\documentclass{ximera}

\input{../../preamble.tex}

\title{Roots and Radicals}

\begin{document}

\begin{abstract}
aspects
\end{abstract}
\maketitle




\begin{definition} \item \textbf{\textcolor{green!50!black}{$n^{th}$ root}}


Let $s$ be a real number. Let $n$ be a natural number.  Then $r$ is an $n^{th}$-root of $s$ provided


\[    r^n = s     \]



The symbol for the $n^{th}$-root of $s$ uses the \textbf{radical sign}, $\sqrt[n]{s}$.

\end{definition}





\subsection*{Even Roots}

This course is the study of the real numbers. As a result we don't take even roots of negative numbers, because a negative number raised to an even power always results in a positive number.



\[   EvenRoot(r) = \sqrt[2n]{r} = r^{\tfrac{1}{2n}}          \]

The square root function is the most popular Core root function


\[   \sqrt{x} = x^{\tfrac{1}{2}}          \]



The domain of the square root function (and all even roots) is $[0, \infty)$, the \textbf{nonnegative} numbers.  






The graph of $y = SR(x) = \sqrt{x}$


\begin{image}
\begin{tikzpicture} 
  \begin{axis}[
            domain=-10:10, ymax=10, xmax=10, ymin=-10, xmin=-10,
            axis lines =center, xlabel=$x$, ylabel=$y$,
            ytick={-10,-8,-6,-4,-2,2,4,6,8,10},
            xtick={-10,-8,-6,-4,-2,2,4,6,8,10},
            ticklabel style={font=\scriptsize},
            every axis y label/.style={at=(current axis.above origin),anchor=south},
            every axis x label/.style={at=(current axis.right of origin),anchor=west},
            axis on top
          ]
          
            %\addplot [line width=2, penColor, smooth, domain=(-9:0),<->] {(x-1)/((x+3)*(x-4))};
            \addplot [line width=2, penColor, smooth, samples=200,domain=(0:9),->] {sqrt(x)};
   
      \addplot[color=penColor,fill=penColor,only marks,mark=*] coordinates{(0,0)};

           

  \end{axis}
\end{tikzpicture}
\end{image}



The square root function begins where the inside of the radical equals $0$.  It then moves in the direction that keeps the inside positive.


All even root/radical functions can be viewed as the Core Square Root function composed with linear function. One linear function affects the inside of the radical sign and the other affects the outside.


\begin{idea} \textbf{\textcolor{blue!55!black}{Square Root Functions}} 


Let $SQ(x) = \sqrt{k}$ be the Core Square Root function.

Let $L_{in}(t) = B \, t + C$ be a linear function.

Let $L_{out}(y) = A \, y + D$ be a linear function.


In general, square root (radical) functions are constructed by composing these component functions.


\[
(L_{out} \circ SQ \circ L_{in})(x) = A \sqrt{B \,x + C} + D
\]


\end{idea}



\begin{summary} \textbf{\textcolor{red!70!yellow}{Analysis}}


Consider the general square root function.

\[  A \sqrt{B \,x + C} + D  \]


\textbf{Domain:} The domain is all real numbers that make the inside nonnegative


\[ B \,x + C >= 0 \]


\textbf{Zeros:} Square root functions can have a single zero or no zeros.


\[  A \sqrt{B \,x + C} + D  = 0 \]

Since square roots of real numbers are always positive (or $0$), this equation has a solution only if $A$ and $D$ have opposite signs.


\textbf{Continuity:} Square root functions are contuous functions.  They are continuous on their domain.



\textbf{End-Behavior:} Square root functions have end-behavior only in the direction that makes the inside positive. In that direction, the value of the function approaches $-\infty$ or $\infty$, which is dictated by the sign of the leading coeffcient.



\textbf{Behavior:} Square root functions are either increasing functions or they are decreasing functions.  

The Chain Rule will give us this information from the behavior of its component functions.


\[
(L_{out} \circ SQ \circ L_{in})(x) = A \sqrt{B \,x + C} + D
\]


$SQ$ is the Core square root function and is an increasing function.  $L_{out}$ and $L_{in}$ are linear functions and their behavior is given by the sign of their leading coefficients.

Then apply the Chain Rule.



\textbf{Global Maximum and Minimum:}  Square root functions have either a global minimum or a global maximum, but not both. This global extreme value occurs at the real number that makes the inside of the radical equal to $0$.


\textbf{Local Maximums and Minimums:} Same as global.


\textbf{Range:} The range is either of the form $[m, \infty)$ or $(-\infty, M]$.

\end{summary}








\begin{example} Square Root Function

The graph of $y = f(v) = \frac{1}{3} \sqrt{5-3v}-4$


\begin{image}
\begin{tikzpicture} 
  \begin{axis}[
            domain=-10:10, ymax=10, xmax=10, ymin=-10, xmin=-10,
            axis lines =center, xlabel=$v$, ylabel=$y$, 
            ytick={-10,-8,-6,-4,-2,2,4,6,8,10},
            xtick={-10,-8,-6,-4,-2,2,4,6,8,10},
            ticklabel style={font=\scriptsize},
            every axis y label/.style={at=(current axis.above origin),anchor=south},
            every axis x label/.style={at=(current axis.right of origin),anchor=west},
            axis on top
          ]
          
          	%\addplot [line width=2, penColor, smooth, domain=(-9:0),<->] {(x-1)/((x+3)*(x-4))};
          	\addplot [line width=2, penColor, smooth, samples=200,domain=(-9:2.67),<-] {0.33*sqrt(5-3*x)-4};
   
 			\addplot[color=penColor,fill=penColor,only marks,mark=*] coordinates{(1.67,-4)};

           

  \end{axis}
\end{tikzpicture}
\end{image}

The domain of $f$ is the interval that makes the inside of the radical nonnegative.  

\[ \left(-\infty, \frac{5}{3} \right] \]


$f$ is the composition of

\begin{itemize}
\item $SQ(x) = \sqrt{k}$ 

\item $L_{in}(t) = 5 - 3 \, t$, a decreasing linear function

\item $L_{out}(y) = \frac{1}{3} \, y + 4$, an increasing linear function
\end{itemize}



\[ inc \circ inc \circ dec = dec  \]


$f$ is a decreasing square root function, which agrees with the graph.


\end{example}















\begin{example} Even Root

The graph of $y = T(k) = \sqrt{-k+5}$


\begin{image}
\begin{tikzpicture} 
  \begin{axis}[
            domain=-10:10, ymax=10, xmax=10, ymin=-10, xmin=-10,
            axis lines =center, xlabel=$k$, ylabel=$y$, 
            ytick={-10,-8,-6,-4,-2,2,4,6,8,10},
            xtick={-10,-8,-6,-4,-2,2,4,6,8,10},
            ticklabel style={font=\scriptsize},
            every axis y label/.style={at=(current axis.above origin),anchor=south},
            every axis x label/.style={at=(current axis.right of origin),anchor=west},
            axis on top
          ]
          
          	%\addplot [line width=2, penColor, smooth, domain=(-9:0),<->] {(x-1)/((x+3)*(x-4))};
          	\addplot [line width=2, penColor, smooth, samples=200,domain=(-9:5),<-] {sqrt(-x+5)};
   
 			\addplot[color=penColor,fill=penColor,only marks,mark=*] coordinates{(5,0)};

           

  \end{axis}
\end{tikzpicture}
\end{image}


Here the inside $-k+5=0$ equals $0$ when $k=\answer{5}$.  That is the start of the domain.  Then the inside is positive $-k+5>0$, when $k$ \wordChoice{\choice[correct]{$<$} \choice {$>$}}  $5$, which means the graph moves up to the left.  The domain is $(-\infty,5]$.




\end{example}














\subsection*{Odd Roots}

This course is the study of the real numbers. As a result we don't take even roots of negative numbers, however we do have odd roots of negative numbers

Odd roots

\[   Odd(x) = \sqrt[2n+1]{x} = x^{\tfrac{1}{2n+1}}          \]

look a lot like the cube root


\[   Odd(x) = \sqrt[3]{x} = x^{\tfrac{1}{3}}          \]



Their domains include all real numbers: $(-\infty, \infty)$.  Odd root functions increase very slowly over this domain, becoming unbounded.







The graph of $y = CubeRoot(x) = \sqrt[3]{x}$


\begin{image}
\begin{tikzpicture} 
  \begin{axis}[
            domain=-10:10, ymax=10, xmax=10, ymin=-10, xmin=-10,
            axis lines =center, xlabel=$x$, ylabel=$y$, 
            ytick={-10,-8,-6,-4,-2,2,4,6,8,10},
            xtick={-10,-8,-6,-4,-2,2,4,6,8,10},
            ticklabel style={font=\scriptsize},
            every axis y label/.style={at=(current axis.above origin),anchor=south},
            every axis x label/.style={at=(current axis.right of origin),anchor=west},
            axis on top
          ]
          
          	%\addplot [line width=2, penColor, smooth, domain=(-9:0),<->] {(x-1)/((x+3)*(x-4))};
          	\addplot [line width=2, penColor, smooth, samples=200,domain=(0:9),->] {x^0.33333};
          	\addplot [line width=2, penColor, smooth, samples=200,domain=(-9:0),<-] {-(-x)^0.33333};
   
 			%\addplot[color=penColor,fill=penColor,only marks,mark=*] coordinates{(0,0)};

           

  \end{axis}
\end{tikzpicture}
\end{image}



The cube root function has a vertical tangent line where the inside of the radical equals $0$.




































\subsection*{Odd Roots}

This course is the study of the real numbers. As a result we don't take even roots of negative numbers, because a negative number raised to an even power always results in a positive number. However, an odd root of a negative number is again a negative number.



\[   OddRoot(r) = \sqrt[2n]{r} = r^{\tfrac{1}{2n}}          \]

The cube root function is the most popular Core odd root function


\[   \sqrt[3]{x} = x^{\tfrac{1}{3}}          \]



The domain of the cube root function (and all odd roots) is $(-\infty, \infty)$.  




The graph of $y = CR(x) = \sqrt[3]{x}$


\begin{image}
\begin{tikzpicture} 
  \begin{axis}[
            domain=-10:10, ymax=10, xmax=10, ymin=-10, xmin=-10,
            axis lines =center, xlabel=$x$, ylabel=$y$,
            ytick={-10,-8,-6,-4,-2,2,4,6,8,10},
            xtick={-10,-8,-6,-4,-2,2,4,6,8,10},
            ticklabel style={font=\scriptsize},
            every axis y label/.style={at=(current axis.above origin),anchor=south},
            every axis x label/.style={at=(current axis.right of origin),anchor=west},
            axis on top
          ]
          
            %\addplot [line width=2, penColor, smooth, domain=(-9:0),<->] {(x-1)/((x+3)*(x-4))};
            \addplot [line width=2, penColor, smooth, samples=200,domain=(-9:0),<->] {-((-x)^(0.333))};

            \addplot [line width=2, penColor, smooth, samples=200,domain=(0:9),<->] {(x)^(0.333)};
   
      %\addplot[color=penColor,fill=penColor,only marks,mark=*] coordinates{(0,0)};

           

  \end{axis}
\end{tikzpicture}
\end{image}






All odd root/radical functions can be viewed as the Core Cube Root function composed with linear function. One linear function affects the inside of the radical sign and the other affects the outside.


\begin{idea} \textbf{\textcolor{blue!55!black}{Cube Root Functions}} 


Let $CR(x) = \sqrt[3]{k}$ an increasing function.

Let $L_{in}(t) = B \, t + C$ be a linear function.

Let $L_{out}(y) = A \, y + D$ be a linear function.


In general, cube root (radical) functions are constructed by composing these component functions.


\[
(L_{out} \circ CR \circ L_{in})(x) = A \sqrt[3]{B \,x + C} + D
\]


\end{idea}



\begin{summary} \textbf{\textcolor{red!70!yellow}{Analysis}}


Consider the general cube root function.

\[  A \sqrt[3]{B \,x + C} + D  \]


\textbf{Domain:} The domain is all real numbers.


\textbf{Zeros:} Cube root functions always have a single.


\[  A \sqrt[3]{B \,x + C} + D  = 0 \]




\textbf{Continuity:} CUbe root functions are contuous functions.   



\textbf{End-Behavior:} Cube root functions have opposite end-behavior on the two sides.  This is dictated by the behavior of the function.



\textbf{Behavior:} Cube root functions are either increasing functions or they are decreasing functions.  

The Chain Rule will give us this information from the behavior of its component functions.


\[
(L_{out} \circ CR \circ L_{in})(x) = A \sqrt[3]{B \,x + C} + D
\]


$CR$ is the Core cube root function and is an increasing function.  $L_{out}$ and $L_{in}$ are linear functions and their behavior is given by the sign of their leading coefficients.

Then apply the Chain Rule.



\textbf{Global Maximum and Minimum:}  Cube root functions do not have a global maximum or minimum.


\textbf{Local Maximums and Minimums:} Cube root functions do not have a local maximums or minimums.


\textbf{Range:} The range is always $(-\infty, \infty)$.

\end{summary}



















\begin{example} Cube Root Function


The graph of $y = g(t) = -4 \sqrt[3]{2t+3}+1$


\begin{image}
\begin{tikzpicture} 
  \begin{axis}[
            domain=-10:10, ymax=10, xmax=10, ymin=-10, xmin=-10,
            axis lines =center, xlabel=$t$, ylabel=$y$, 
            ytick={-10,-8,-6,-4,-2,2,4,6,8,10},
            xtick={-10,-8,-6,-4,-2,2,4,6,8,10},
            ticklabel style={font=\scriptsize},
            every axis y label/.style={at=(current axis.above origin),anchor=south},
            every axis x label/.style={at=(current axis.right of origin),anchor=west},
            axis on top
          ]
          
          	\addplot [line width=2, penColor, smooth, samples=200,domain=(-1.49:8),->] {-4*(2*x+3)^0.33333+1};

          	\addplot [line width=2, penColor, smooth, samples=200,domain=(-6:-1.49),<-] {4*(-2*x-3)^0.33333+1};

          	%\addplot[color=penColor,fill=penColor,only marks,mark=*] coordinates{(-3,0)};

           

  \end{axis}
\end{tikzpicture}
\end{image}


Here the inside $2t+3$ equals $0$ when $v=\answer{\frac{-2}{3}}$.  The graph has a vertical tangent line at $\left( -\frac{2}{3}, 1 \right)$.  




$g$ is the composition of

\begin{itemize}
\item $CR(x) = \sqrt[3]{k}$, an increasing function

\item $L_{in}(t) = 2t+3$, an increasing linear function

\item $L_{out}(y) = -4y + 1$, a decreasing linear function
\end{itemize}



\[ dec \circ inc \circ inc = dec  \]


$g$ is a decreasing cube root function, which agrees with the graph.


\end{example}
































\begin{center}
\textbf{\textcolor{green!50!black}{ooooo-=-=-=-ooOoo-=-=-=-ooooo}} \\

more examples can be found by following this link\\ \link[More Examples of Analysis]{https://ximera.osu.edu/csccmathematics/precalculus1/precalculus1/libraryAnalysis1/examples/exampleList}

\end{center}






\end{document}
